\section{Deployment}

\subsection{Bygging, lagring og lasting}

Sluttproduktet er en lokal Streamlit-demonstrasjon, ikke en offentlig driftet
tjeneste. Kommandoen \texttt{make build-app-model} kjører tre tydelige steg:
nedlasting av trafikkdata, normalisering med bevarte målegap og trening av den
frosne kandidaten. Treningsskriptet gjenskaper CV-rapporten og sluttevalueringen,
trener modellen på alle gyldige rader før 1. januar 2026 og lagrer en modellpakke
med Joblib i \texttt{models/traffic\_hgb\_latency\_robust\_v1.joblib}. Pakken
inneholder den trente HGB-en, den ordnede listen over \emph{features} og metadata om
modellversjon, lagpakke, utvalgsperiode, treningsperiode og datakilde. Den tilhørende
JSON-filen viser 34 159 treningsrader.

Joblib gir \texttt{dump} og \texttt{load} for serialisering av Python-objekter
\cite{joblib-persistence}. Artefakten er ignorert av Git og må bygges lokalt. Appen
laster bare den forventede lokale filen og gir en konkret byggekommando dersom den
mangler eller ikke kan leses. Fordi Joblib bygger på pickle, bør en operativ løsning
aldri laste en modellfil fra en uklar eller ukontrollert kilde.

\subsection{Brukerflyt og prognose}

Appen startes med \texttt{streamlit run app/streamlit\_app.py}. Streamlit er et
Python-rammeverk for interaktive dataapplikasjoner \cite{streamlit-docs}. Brukeren
velger enten \enquote{Neste time} eller et historisk tidspunkt og trykker
\enquote{Vis prognose}. Historiske valg begrenses til reelle lokale klokketimer:
begge forekomster av en gjentatt time ved vintertid kan velges, mens en
ikke-eksisterende time ved sommertid ikke tilbys. Prognoseintervallet vises med
UTC-offset og \texttt{Europe/Oslo}.

Ved hvert klikk lastes modellen, og appen henter nødvendig historikk fra Traffic Data
API. Tidspunktet som skal predikeres, er $t+1$, mens $t$ er ankertidspunktet for
modellens \emph{features}. Koden bygger kalenderverdier og slår opp trafikk ved $t-24$,
$t-48$ og $t-168$. Hvis en av disse verdiene mangler, blir det ingen prediksjon.
Hvis alle finnes, viser appen forventet biltrafikk i kjøretøy per time, en graf med
de siste 72 timene og en egen diamant for prognosen. Målegap og upubliserte timer
vises som åpne hull, uten linjer som skjuler manglende data.

Appens sammenligning med forrige uke må ikke forveksles med modellens ukeslag.
Grensesnittet viser observert trafikk ved \emph{target minus 168 timer}. Siden
\emph{target} er $t+1$, tilsvarer dette $t-167$. Det er samme tidsposisjon som den
operative forrige-uke-baselinen bruker. Modellens \emph{feature} ved $t-168$ er derimot
knyttet til ankertidspunktet $t$ og ligger 169 timer før targetverdien. Verdiene har
derfor ulike tidspunkter, selv om begge omtales som ukentlig historikk. Ved sommertid
og vintertid kan også det lokale klokkeslettet forskyves når forskjellen beregnes som
eksakt forløpt tid.

Sammenligningsverdien er valgfri: mangler den, kan en gyldig modellprognose fortsatt
vises. Prosentforskjellen beregnes bare når sammenligningen er positiv. En utfellbar
seksjon viser modellnavn, \emph{features}, evalueringsresultat og de konkrete inndataene.
Datastatus viser kilde, hentetidspunkt og siste tilgjengelige observasjon i uttrekket.

\subsection{Feilhåndtering, demonstrasjon og drift}

Appen skiller mellom manglende eller skadet modellartefakt, feil fra API-et og
utilstrekkelig trafikkhistorikk. Brukeren får et kort neste steg, mens tekniske
detaljer om manglende inndata kan åpnes ved behov. Hvert klikk gjør et nytt API-kall;
det brukes ikke mellomlager eller lokal database i denne versjonen. Historiske prognoser
bruker data som er tilgjengelige nå. De rekonstruerer ikke hva API-et hadde
publisert på det historiske prognosetidspunktet, og prognoser innen treningsperioden
er demonstrasjoner, ikke nye testresultater.

En videodemonstrasjon er lenket fra prosjektets README og oppfyller oppgavens krav
om å vise nettsiden uten offentlig publisering:
\href{https://github.com/user-attachments/assets/5598cc2d-ddde-4813-93b8-5eb32f29f990}{Bergen
Traffic Predictor videodemonstrasjon}. Løsningen kan kjøres på en vanlig lokal
datamaskin med Python 3.11 eller nyere. Den krever ingen Frost-nøkkel fordi vær ikke
inngår i sluttproduktet.

Før reell drift må systemet få vedvarende logging og overvåking. Minstekrav er
API-feilrate, andel prognoser som mangler ett eller flere lag, alder på siste
observasjon og endringer i fordelingen av trafikkvolum. Når fasiten senere blir
observert, bør MAE og RMSE følges over tid og fordeles på time og ukedag. Vedvarende
avvik kan skyldes datadrift, sensorendring, vegomlegging eller endret reisemønster.
Modellen bør trenes på nytt etter en fast vurderingsplan eller ved dokumentert drift,
men en ny versjon må velges på kronologisk validering og måles på en ny, urørt
periode. API-tilgjengelighet, datalisens, artefaktintegritet og en manuell
reserveprognose må også inngå i en produksjonsplan.

Overvåkingen bør skille datafeil fra modelldrift. Manglende lag eller en plutselig
økning i API-feil peker mot datakjeden, mens stabile inndata med økende
prediksjonsfeil kan tyde på modelldrift. Modell- og dataversjon bør logges sammen med
hver prognose slik at en feil kan spores. Før automatisk ny trening innføres, må det
finnes en godkjenningsmekanisme som hindrer at et midlertidig målegap eller en
unormal periode blir gjort til ny normal.
